\chapter*{Notation}
\label{frontmatter:notation}
\addcontentsline{toc}{chapter}{Notation}

A consolidated table of the symbols and conventions used across all twelve volumes. Per-volume notation that does not generalise is introduced where it is first needed. The cross-volume vocabulary is collected here.

\section*{Quantities and sets}

\begin{tabular}{ll}
  \toprule
  Symbol & Meaning \\
  \midrule
  $\R, \N, \Z, \Q, \C$ & real, natural, integer, rational, complex numbers \\
  $\F$ & a generic field \\
  $\set{\,\cdot\,}$ & a set \\
  $\inner{\,\cdot\,,\,\cdot\,}$ & inner product \\
  $\norm{\,\cdot\,}$ & norm (Euclidean unless stated) \\
  $\abs{\,\cdot\,}$ & absolute value or cardinality \\
  $\oom{f(n)}$ & order of magnitude (asymptotic) \\
  \bottomrule
\end{tabular}

\section*{Vectors, matrices, tensors}

\begin{tabular}{ll}
  \toprule
  Symbol & Meaning \\
  \midrule
  $\vect{x}, \vect{F}, \vect{u}$ & vectors (bold) \\
  $\mat{A}, \mat{B}$ & matrices (bold uppercase) \\
  $\tens{T}$ & tensor (bold sans-serif) \\
  $\unitvec{x}, \unitvec{n}$ & unit vectors \\
  $\mat{A}\transpose, \mat{A}\inv$ & transpose and inverse \\
  \bottomrule
\end{tabular}

\section*{Calculus and operators}

\begin{tabular}{ll}
  \toprule
  Symbol & Meaning \\
  \midrule
  $\dd x$ & differential element of $x$ \\
  $\dv{f}{x}$, $\pdv{f}{x}$ & ordinary and partial derivative \\
  $\grad f$ & gradient \\
  $\divg \vect{F}$ & divergence \\
  $\curl \vect{F}$ & curl \\
  $\laplacian f$ & Laplacian \\
  \bottomrule
\end{tabular}

\section*{Probability and statistics}

\begin{tabular}{ll}
  \toprule
  Symbol & Meaning \\
  \midrule
  $\Pr(A)$ & probability of event $A$ \\
  $\E[X]$ & expectation of random variable $X$ \\
  $\Var(X), \Cov(X,Y), \Cor(X,Y)$ & variance, covariance, correlation \\
  $\Hent(X)$ & Shannon entropy \\
  $\Imut(X;Y)$ & mutual information \\
  $\KL(P \| Q)$ & Kullback-Leibler divergence \\
  \bottomrule
\end{tabular}

\section*{Units}

The book uses SI throughout, except where a non-SI unit is the working unit of a discipline (electron-volts in Volume IV, parsecs in Volume XII discussions of planetary engineering, gauge in Volume VIII when reproducing US-standard data sheets). All non-SI usage is annotated and the SI equivalent is given in parentheses on first use in any chapter.

The seven SI base units (\si{\metre}, \si{\kilogram}, \si{\second}, \si{\ampere}, \si{\kelvin}, \si{\mole}, \si{\candela}) and the standard prefixes are taken as known by the reader. Volume I introduces them rigorously for the reader who is not yet at home with them.

\section*{Editorial markers}

\begin{tabular}{ll}
  \toprule
  Marker & Meaning \\
  \midrule
  \TODO{} & editorial placeholder, removed before publication \\
  \engterm{term} & technical term in bold on first definition \\
  \engdef{phrase} & emphasised phrase in italic \\
  \bottomrule
\end{tabular}

\noindent
\TODO{notation table is the cross-volume consolidated index. Per-volume notation is introduced where it is first needed; the per-volume notation tables in each volume's opening chapter are populated as those chapters are written.}
